\chapter{Interbank Matching Mechanisms}

The position of the matching stage within the daily simulation process
is described in Section~3.1. This chapter focuses only on how eligible
lenders and borrowers are connected, how transaction amounts and rates
are determined, and how the resulting trades differ between centralized
matching and decentralized RFQ. Clearing, market updates, bank-state
transitions, and systemic-risk recording are not repeated here.

\section{Centralized Matching Mechanism}

Centralized matching is treated as a high-coordination benchmark. Bank
roles and trading intentions are updated on each business day and the
current orders enter a market-wide allocation routine on that same day.
The baseline cycle length is therefore one business day. A longer cycle
can be imposed as a structural robustness exercise, but it is not used in
the principal Centralized--RFQ comparison.

The operational objective is to provide a deterministic,
lender-return-oriented benchmark under full eligible-market information.
One coordinating principal gives priority to borrowers associated with
the highest feasible midpoint settlement rate and then considers their
feasible lenders in descending midpoint-rate order. Conditional on that
priority, it attempts to assemble each borrower's complete demand and
commits the provisional allocation only when full funding is feasible.
The mechanism therefore does not minimize borrower financing costs and is
not presented as a global welfare or maximum-cardinality optimizer. Its
purpose is to apply one transparent common rule to all eligible orders.

In each period, eligible lending supply and borrowing demand are
considered jointly across the market. The mechanism determines feasible
lender--borrower pairs and allocates funds subject to interest-rate
compatibility, lender capacity, borrower demand, the single-loan cap,
and network-degree restrictions.

Let \(\mathcal{L}_{t}\) denote the set of active lenders and let
\(\mathcal{B}_{t}\) denote the set of active borrowers in period \(t\).
The available supply of lender \(i\) is denoted
by \(S_{i}^{(t)}\), while the funding demand of borrower \(j\) is denoted
by \(D_{j}^{(t)}\). These quantities are recalculated from the current
bank states before daily allocation.

Throughout Chapter~4, \(x_{ij}^{(t)}\) denotes a new loan from lender
\(i\) to borrower \(j\) in period \(t\). The first index therefore
identifies the lender, while the second identifies the borrower. This
orientation differs from the due-payment liability matrix in Section~5.1,
where the first index identifies the debtor.

A lender--borrower pair is eligible only when \(i \in \mathcal{L}_{t}\),
\(j \in \mathcal{B}_{t}\), \(i \neq j\), both banks remain active, and
the applicable degree and rate constraints are satisfied. Rollover
status does not remove a bank from ordinary matching in the current
implementation.

Let \(r_{i,\min}^{(t)}\) denote the minimum lending rate accepted by
lender \(i\), and let \(r_{j,\max}^{(t)}\) denote the maximum borrowing
rate accepted by borrower \(j\). The interest-rate surplus of a
lender--borrower pair is defined as:

\begin{equation}
\Delta r_{ij}^{(t)}=r_{j,\max}^{(t)}-r_{i,\min}^{(t)}\label{eq:rate-surplus}.
\end{equation}

A positive interest-rate surplus indicates that the pair is compatible.
The executable routine uses the feasible midpoint settlement rate to
implement the lender-return-oriented priority stated above. Reverse
sorting is therefore intentional: higher feasible midpoint rates are
processed before lower rates, rather than minimizing the borrower's
financing cost. This priority is followed by the full-funding
commit-or-rollback rule.

The centralized allocation is governed by the following feasibility and
funding conditions:

\begin{equation}
\begin{aligned}
r_{i,\min}^{(t)}&\leq r_{j,\max}^{(t)},&&i\in\mathcal L_t,\ j\in\mathcal B_t,\ i\neq j,\\
0\leq x_{ij}^{(t)}&\leq\min\left\{S_{i,\mathrm{rem}}^{(t)},D_{j,\mathrm{rem}}^{(t)},\bar x\right\},&&i\in\mathcal L_t,\ j\in\mathcal B_t,\\
\sum_{j\in\mathcal B_t}x_{ij}^{(t)}&\leq S_i^{(t)},&&i\in\mathcal L_t,\\
\sum_{i\in\mathcal L_t}x_{ij}^{(t)}&=z_j^{(t)}D_j^{(t)},&&j\in\mathcal B_t,\\
z_j^{(t)}&\in\{0,1\},&&j\in\mathcal B_t.
\end{aligned}\label{eq:centralized-conditions}
\end{equation}
\vspace{-0.4\baselineskip}
\noindent Here, \(\bar x\) denotes the upper limit on a single loan, while
\(z_{j}^{(t)}\) indicates whether borrower \(j\) is fully funded. When
\(z_{j}^{(t)} = 1\), the borrower's complete demand is covered and the
tentative transactions are committed. When \(z_{j}^{(t)} = 0\), complete
funding cannot be assembled, the tentative allocations are cancelled,
and the borrower receives no new centralized loan in that period.

The implementation applies this condition through a ranking-and-rollback
procedure rather than by solving a separate mathematical programming
problem. Feasible pairs are ranked, and lender supply is tentatively
allocated subject to the single-loan and degree limits. The allocation
is committed only when the borrower's full demand can be covered;
otherwise, the tentative transactions are rolled back and the lenders'
available supplies are restored.

Committed trades are processed through the same balance-sheet accounting,
contract recording, and clearing modules used for RFQ trades. However,
the two mechanisms should not be interpreted as differing only in
counterparty access. Both operate daily, but centralized matching combines
market-wide visibility and an all-or-nothing borrower rule, whereas RFQ
combines local discovery, bilateral quotation, risk screening, and retained
partial funding.

\section{Decentralized RFQ Matching Mechanism}

Decentralized request-for-quote matching is executed on each business
day after active banks have been assigned lender, borrower, or
non-participant roles. This section focuses on the
bilateral enquiry, quotation, and transaction-acceptance procedures used
during the daily matching stage.

Unlike the centralized principal, RFQ has no market-wide coordinator.
Each borrower initiates search, observes only a limited local lender pool,
and requests bilateral quotations. Lenders determine whether and on what
terms to quote according to their remaining funds, reservation rates, and
the borrower's condition. A borrower may conduct several enquiry rounds
when its demand is not fully satisfied and retains acceptable partial
trades already completed.

The RFQ mechanism begins with the formation of trading intentions. In
the implementation, each intention records the bank identifier, market
role, relevant reservation rate, and intended transaction quantity. A
lender specifies its minimum acceptable lending rate and available
supply, whereas a borrower specifies its maximum acceptable borrowing
rate and desired funding amount. The central bank, inactive banks, and
banks assigned as non-participants are excluded from ordinary RFQ
matching.

For lender \(i\), loanable supply is determined by the amount of liquid
assets remaining above its target liquidity buffer. Let
\(LA_{i}^{(t)}\) denote the liquid assets of lender \(i\) in period
\(t\), and let \(LB_{i}^{(t)}\) denote its target liquidity buffer. The
available lending supply is defined as:


\begin{equation}
S_i^{(t)}=\max\left\{0,\,\mathrm{LA}_i^{(t)}-\mathrm{LB}_i^{(t)}\right\}\label{eq:rfq-supply}.
\end{equation}


The target liquidity buffer reflects the bank's current liabilities,
expected funding outflows, reserve requirements, and liquidity target.
This constraint prevents a lender from supplying funds that are needed
to maintain its own liquidity position.

Borrowing demand may arise from either a liquidity shortage or an
investment opportunity. For borrower \(j\), the liquidity-gap component
is defined as:


\begin{equation}
D_{j,\mathrm{gap}}^{(t)}=\max\left\{0,\,\mathrm{LB}_j^{(t)}-\mathrm{LA}_j^{(t)}\right\}\label{eq:rfq-gap-demand}.
\end{equation}
\vspace{-0.4\baselineskip}
\noindent Here, \(LA_{j}^{(t)}\) denotes the borrower's liquid assets, while
\(LB_{j}^{(t)}\) denotes its target liquidity buffer. A positive value
indicates that the borrower's available liquidity is below the target
level.

The borrower may also demand additional funds when expected project
returns justify expansion. Total borrowing demand is therefore defined
as:


\begin{equation}
D_j^{(t)}=D_{j,\mathrm{gap}}^{(t)}+D_{j,\mathrm{opp}}^{(t)}\label{eq:rfq-total-demand}.
\end{equation}
\vspace{-0.4\baselineskip}
\noindent where \(D_{j,opp}^{(t)}\) represents opportunity-driven demand. This
component depends on the expected project return, the cost of interbank
borrowing, the borrower's risk appetite, and the current market
environment. Consequently, a bank may enter the RFQ market either
because it faces a liquidity shortage or because it identifies a
profitable investment opportunity.

The RFQ market operates through a limited number of bilateral enquiry
rounds. At initialization, commercial banks are placed on a randomly
permuted ring and each bank is acquainted with its two nearest neighbors
on either side. These information links permit enquiries but do not create
loans or balance-sheet exposures. In each period a borrower constructs one
local lender pool in the following order: lenders on outstanding contracts,
historical counterparties, initial acquaintances and current one-hop
exposure neighbors, two-hop referrals, and at most one randomly explored
lender. Outstanding lenders do not consume the new-discovery budget; at
most eight other lenders may enter the pool.

The borrower may conduct up to four RFQ rounds and contact at most four
lenders from the same local pool in each round. A node-level degree cap of
12 applies separately to both lender and borrower. Existing links remain
eligible, while a new link requires both nodes to be below the cap.
Degree-blocked lenders and lenders with no bilateral room are filtered
before the four candidate positions are filled, so they do not repeatedly
occupy a position in later rounds. A lender that can supply only part of
the requested amount remains eligible only while usable cash and bilateral
room remain. These rules improve execution within the local information set
without giving RFQ a market-wide view.

Each locally feasible pair is evaluated with the version-6 GNN matcher and
a teacher-aligned risk score. Ranking gives priority to low borrower
origination risk and penalizes post-trade bilateral exposure, while GNN
compatibility, relationship information, and the rate margin act as
secondary terms. The GNN therefore changes selection within the local pool;
it does not expand that pool.

Each contacted lender determines whether to submit a quotation. For
lender \(i\) and borrower \(j\), the quoted interest rate is defined as:


\begin{equation}
q_{ij}^{(t)}=r_{i,\min}^{(t)}+\mu_{ij}^{(t)}\label{eq:rfq-quote}.
\end{equation}
\vspace{-0.4\baselineskip}
\noindent where \(r_{i,\min}^{(t)}\) denotes lender \(i\)'s minimum acceptable
lending rate, and \(\mu_{ij}^{(t)}\) denotes the borrower-specific risk
markup. The markup reflects the lender's assessment of the borrower's
capital, liquidity, solvency, and general financial condition. A
financially weaker borrower may therefore receive a higher quoted rate.

A quotation is acceptable only when the quoted rate does not exceed the
borrower's maximum acceptable borrowing rate:


\begin{equation}
q_{ij}^{(t)}\leq r_{j,\max}^{(t)}.
\end{equation}
\vspace{-0.4\baselineskip}
\noindent where \(r_{j,\max}^{(t)}\) denotes the highest borrowing rate that
borrower \(j\) is willing to accept. Quotations that exceed this limit
are rejected.

Among the acceptable quotations, the borrower gives priority to offers
with favorable risk-adjusted scores. The amount of an accepted
transaction is constrained by the lender's remaining supply, the
borrower's remaining demand, the borrower's balance-sheet borrowing room,
and a bilateral gross-exposure limit. Let \(x_{ij}^{(t)}\) denote the
amount lent by lender \(i\) to borrower \(j\). If
\(E_{ij,t}^{\mathrm{gross}}\) is their existing gross exposure, the
pair-specific room is

\begin{equation}
\bar x_{ij,t}^{\mathrm{pair}}
=\min\left\{
\max\{0,0.25C_{i,t}-E_{ij,t}^{\mathrm{gross}}\},
0.5D_{j,t}^{\mathrm{initial}},\bar x
\right\}.
\label{eq:rfq-pair-cap}
\end{equation}

The transaction amount is then bounded by


\begin{equation}
x_{ij}^{(t)}=\min\left\{S_{i,\mathrm{rem}}^{(t)},D_{j,\mathrm{rem}}^{(t)},
R_{j,t}^{\mathrm{borrow}},\bar x_{ij,t}^{\mathrm{pair}}\right\}
\label{eq:rfq-amount}.
\end{equation}
\vspace{-0.4\baselineskip}
\noindent where \(S_{i,rem}^{(t)}\) is lender \(i\)'s remaining supply,
\(D_{j,rem}^{(t)}\) is borrower \(j\)'s remaining demand, and
\(R_{j,t}^{\mathrm{borrow}}\) is its remaining balance-sheet borrowing room.
The bilateral limit is applied to gross principal plus arrears and is not
netted against claims in the opposite direction.

A transaction is completed only when the calculated amount is no smaller
than the minimum trade size \(m\):


\begin{equation}
x_{ij}^{(t)}\geq m.
\end{equation}


Under the baseline setting, \(m=10^{-9}\). This floor is retained in the
code as a numerical filter against degenerate zero-size trades, but it is
effectively inactive for economically relevant amounts and does not screen
small network links in the reported experiments.

After a transaction is accepted, the remaining supply of lender \(i\) is
updated as:


\begin{equation}
S_{i,\mathrm{rem}}^{(t)}\leftarrow S_{i,\mathrm{rem}}^{(t)}-x_{ij}^{(t)}.
\end{equation}


and the remaining demand of borrower \(j\) is updated as:


\begin{equation}
D_{j,\mathrm{rem}}^{(t)}\leftarrow D_{j,\mathrm{rem}}^{(t)}-x_{ij}^{(t)}.
\end{equation}


The borrower continues to compare acceptable quotations until its demand
is satisfied, the available quotations are exhausted, or the maximum
number of RFQ rounds is reached.

Unlike the centralized mechanism, RFQ permits partial funding.
Transactions that have already been accepted are retained even when the
borrower's total demand cannot be fully covered. A borrower may
therefore complete several bilateral loans while still having positive
unmet demand at the end of the matching stage.

Accepted trades are processed according to the common
contract-accounting and timing rules described in Section~3.1. Each
completed trade specifies the lender, borrower, principal amount,
interest rate, origination date, and maturity conditions. The detailed
settlement and clearing treatment of these contracts is presented in
Chapters 4.4 and 5.

The RFQ procedure records several indicators describing its immediate
funding outcomes. Total transaction volume in period \(t\) is defined as:


\begin{equation}
V_t=\sum_i\sum_j x_{ij}^{(t)}\label{eq:trading-volume}.
\end{equation}


The funding satisfaction ratio is defined as:


\begin{equation}
\mathrm{FS}_t=\frac{V_t}{\max\left\{D_t,\varepsilon\right\}}\label{eq:funding-satisfaction}.
\end{equation}
\vspace{-0.4\baselineskip}
\noindent where \(D_{t}\) denotes total borrowing demand in period \(t\), and
\(\varepsilon\) is a small positive value used to avoid division by zero.

Aggregate unmet demand is defined as:


\begin{equation}
\mathrm{UD}_t=\max\left\{0,\,D_t-V_t\right\}\label{eq:unmet-demand}.
\end{equation}


The volume-weighted average transaction rate is calculated as:


\begin{equation}
\bar q_t=\frac{\displaystyle\sum_i\sum_j q_{ij}^{(t)}x_{ij}^{(t)}}{\max\left\{\displaystyle\sum_i\sum_j x_{ij}^{(t)},\varepsilon\right\}}\label{eq:average-quote}.
\end{equation}


These indicators describe the direct outcome of the RFQ matching stage.
They should be distinguished from clearing defaults, bank failures, and
capital stress, which are evaluated when the corresponding contractual
payments become due in later periods. Section~4.3 compares these
characteristics with those of the centralized matching mechanism.

\section{Comparison between Centralized and Decentralized Matching}

The implemented centralized mechanism and implemented RFQ mechanism use the same bank population,
balance-sheet rules, regulatory constraints, contract accounting, and
due-payment clearing procedure. Their transaction-formation rules,
however, differ along both a cross-sectional and a temporal dimension.
Centralized matching executes a market-wide allocation every business day
and applies an all-or-nothing funding rule at the borrower level. RFQ is
also executed every business day,
limits each borrower to local lender enquiries, and retains acceptable
partial trades even when total demand is not fully covered.

The main differences are summarized below.

\begin{table}[htbp]
\centering
\caption{Comparison of centralized matching and decentralized RFQ}
\label{tab:matching-comparison}
\small
\renewcommand{\arraystretch}{1.18}
\begin{tabularx}{\textwidth}{@{}>{\raggedright\arraybackslash}p{0.20\textwidth}XX@{}}
\toprule
\textbf{Aspect} & \textbf{Centralized matching} & \textbf{Decentralized RFQ}\\
\midrule
Decision scope & Market-wide allocation among eligible banks & Local enquiry and bilateral selection\\
Trading frequency & Market-wide allocation every business day & Matching and quotation every business day\\
Counterparty access & All feasible lenders considered jointly & One local pool: at most eight newly discovered lenders plus outstanding lenders; at most four contacted per round\\
Funding rule & All-or-nothing at the borrower level & Accepted partial funding may be retained\\
Pricing & Reservation rates and coordinated ranking & Bilateral quotes and borrower-specific risk markups\\
Unmet demand & Failure to assemble complete funding in the current period & Limited local capacity, bilateral caps, rejected quotes, or insufficient supply\\
Network formation & Daily coordinated formation of market-wide exposures & Repeated local enquiries and accepted bilateral quotations\\
Potential risk channel & Broad access may concentrate allocation in attractive counterparties & Local funding gaps, repeated counterparties, or residual bilateral concentration\\
\bottomrule
\end{tabularx}
\end{table}

The centralized design can combine funds from several lenders to satisfy
one borrower, but a borrower receives no funding when complete demand
cannot be assembled. Counterparty access is broad, so its principal
friction is the all-or-nothing completion rule rather than infrequent
market access.

RFQ produces a different funding pattern. A borrower can seek funding on
each business day, but it observes only a limited subset of lenders in
each enquiry round. Individually accepted trades are retained even when
total demand remains partly unsatisfied. Daily access can therefore
relieve some immediate funding pressure, while limited search and
bilateral pricing can still leave positive unmet demand.

These rules also influence network formation. Daily centralized allocation
can direct large amounts toward a relatively small set of attractive
counterparties. Daily RFQ can repeatedly create and renew bilateral
relationships, so a decentralized network need not be sparser than the
centralized one. The number of links, exposure concentration, repeated
counterparties, and component size are therefore treated as empirical
simulation outcomes rather than imposed properties of either mechanism.

The implications for systemic risk are ambiguous before the simulations
are run. The implemented centralized mechanism combines full eligible-order
visibility, descending rate priority, and all-or-nothing borrower funding.
The implemented RFQ mechanism combines local discovery, bilateral pricing
and screening, and retained partial funding. Observed differences therefore
reflect these bundled features and must not be attributed solely to
centralization or decentralization.

The comparison is evaluated using four groups of outcomes. Funding
performance is measured through actual lending, funding satisfaction,
and unmet demand. Network formation is examined through the number and
distribution of links, exposure concentration, and repeated trading
relationships. Settlement outcomes include actual payments, payment
shortfalls, and clearing defaults. System-wide effects are assessed
using \(\mathrm{FR}_t\), \(\mathrm{CBS}_t\), \(\mathrm{CGR}_t\), and the composite
indicator \(\mathrm{SR}_t\). Because both mechanisms operate daily, the
estimated difference primarily reflects information scope, allocation and
pricing rules, GNN-assisted local screening, and retained partial funding;
it is not a comparison driven by unequal transaction frequency.

\section{Rollover and Coupon Mechanism}

Interbank contracts persist through a dated contract book that records
the lender, borrower, principal, remaining principal, settlement rate,
creation step, maturity step, repayment type, coupon rate, total tenor,
and the number of completed installment periods. The current rollover
implementation selects the repayment schedule when the trade is
originated. Hence the ON/OFF comparison changes contractual maturity
from the beginning; it does not wait for a bullet loan to mature before
deciding whether to roll its principal.

When rollover is disabled, a new trade is a one-business-day
single-payment contract. If \(P_c\) is principal and
\(r_c^{\mathrm{settle}}\) is the transaction rate, the amount due in the
next settlement cycle is

\begin{equation}
D_{c,t}=P_c\left(1+r_c^{\mathrm{settle}}\right).
\label{eq:single-payment-due}
\end{equation}

Any residual beyond the clearing tolerance after the common
Eisenberg--Noe settlement causes a hard payment default. The unpaid claim
is retained for creditor recovery and write-off accounting rather than
being silently deleted.

When rollover is enabled, the entire notional is recorded immediately as
an amortizing installment contract. Its tenor \(N_c\) lies between 5 and
20 business days and increases with principal: amounts at or below 50
simulation units receive five periods, amounts at or above 2000 receive
20 periods, and intermediate amounts are linearly interpolated and
rounded. No fixed rollover fraction is used.

The coupon rate of the rolled contract is based on the original
settlement rate plus a tenor-dependent rollover spread. Let
\(P_{c,t}^{-}\) denote the remaining principal immediately before an
installment payment in period \(t\), and let
\(r_c^{\mathrm{coupon}}\) denote the coupon rate. The interest due in
that period is

\begin{equation}
I_{c,t}
=
r_c^{\mathrm{coupon}}P_{c,t}^{-}.
\label{eq:coupon_payment}
\end{equation}

Let \(k_{c,t}\) be the number of scheduled installments already processed. The
number of remaining installment periods is

\begin{equation}
R_{c,t}=N_c-k_{c,t}.
\label{eq:remaining_periods}
\end{equation}

The principal amount scheduled for the current installment is

\begin{equation}
A_{c,t}^{\mathrm{inst}}
=
\begin{cases}
\displaystyle \frac{P_{c,t}^{-}}{R_{c,t}}, & R_{c,t}>1,\\[6pt]
P_{c,t}^{-}, & R_{c,t}=1.
\end{cases}
\label{eq:installment_principal}
\end{equation}

Arrears from preceding periods are added to the current coupon and
scheduled principal. Thus the amount submitted to clearing is

\begin{equation}
D_{c,t}=Q_{c,t}^{\mathrm{arrears}}+I_{c,t}+A_{c,t}^{\mathrm{inst}}.
\label{eq:contract_due_payment}
\end{equation}

All contracts due on the same day enter one gross Eisenberg--Noe matrix;
there is no contract-by-contract or interest-before-principal clearing
order. After the settlement, remaining scheduled principal is updated as

\begin{equation}
P_{c,t}^{+}
=
\max\left\{0,\,P_{c,t}^{-}-A_{c,t}^{\mathrm{inst}}\right\},
\label{eq:remaining_principal}
\end{equation}

and the installment counter becomes

\begin{equation}
k_{c,t}^{+}=k_{c,t}+1.
\label{eq:period_update}
\end{equation}

If the payment is incomplete, the unpaid total becomes next-period
arrears. The amortization clock still advances, and any principal included
in the missed installment is preserved in arrears rather than erased.
Consecutive misses are recorded as a delinquency statistic but do not by
themselves cause hard default when rollover is enabled. After the planned
tenor, any residual principal or arrears continues to be requested each
business day until settled or resolved through bank default.

A bank with an active installment remains eligible for ordinary
interbank borrowing. The current rollover functions do not remove such a
bank from the borrower set merely because installment debt or arrears are
outstanding. Rollover therefore acts as a maturity and payment-pressure
buffer rather than as an automatic borrowing prohibition.

For both single-payment and installment contracts, only cash flows scheduled
for the current business day enter Eisenberg--Noe clearing. If
\(I_{c,t}\) denotes current interest and \(A_{c,t}\) denotes current
principal due, future coupon payments and remaining principal are retained in the
contract book and excluded from the current-period due-payment matrix.
The due flows are aggregated into the gross debtor-to-creditor liability
matrix described in Chapter~5, so reciprocal obligations are preserved
separately rather than bilaterally netted before clearing.

The purpose of the rollover mechanism is therefore narrow and explicit.
It does not represent a complete refinancing market. It converts a
new interbank loan from next-day single payment into a 5--20-period
amortizing schedule and thereby spreads repayment pressure across future
business days. This
allows the simulations to test whether changing the timing of principal
payments, especially in combination with policy support, delays or
amplifies the realization of systemic stress.
